\section*{Highpass VCVS Butterworth}
\noindent\colorbox{orange}{
  \begin{minipage}{1.0\linewidth}
    \paragraph{Richiesta (a)}
    Dimostrare, discutendone la risposta ai limiti del dominio delle frequenze reali sulla base
    di semplici argomenti di natura qualitativa, che il circuito proposto svolga effettivamente
    la funzione di filtro passa-alto.
  \end{minipage}
}

\paragraph{Risposta (a)}
\begin{itemize}
\item
  Per $f \to 0$ il condensatore $C_1$ apre il ramo dunque $V_+ = 0$. Se l'opamp
  opera in regime di linearit\'a $V_{OUT} = A_dV_D$
  \begin{equation*}
    V_{OUT} = A_d(V_+ - V_-) = A_d(0 - V_{OUT}) \implies \colorbox{blu}{\boxed{V_{OUT} = \frac{0}{1 + A_d} = 0}}
  \end{equation*}
  nel limite di basse frequenze il circuito ha un andamento compatibile con quello di
  un circuito filtro passa alto.
\item Per $f \to +\infty$ i condensatori $C_1$ e $C_2$ ``corto-circuitano'' i propri rami
  cos\'i che $V_+ = V_s$. Se l'opamp opera in regime di linearit\'a $V_{OUT} = A_dV_D$
  \begin{equation*}
    V_{OUT} = A_d(V_+ - V_-) = A_d(V_s - V_{OUT}) \implies \colorbox{blu}{\boxed{V_{OUT} = \frac{A_d}{1 + A_d}V_s}}
  \end{equation*}
  nel limite di alte frequenze il circuito ha un andamento compatibile con quello di
  un circuito filtro passa alto.
\end{itemize}

\noindent\colorbox{orange}{
  \begin{minipage}{1.0\linewidth}
    \paragraph{Richiesta (b)}
    Nell’ipotesi di idealit\'a dell’operazionale ed assumendo C1 = C2 = C e R1=2, R2=2R,
    derivare analiticamente un’espressione della trasformata di Laplace della funzione di
    trasferimento $A_v(s)=V_{OUT}(s)/V_S(s)$.
  \end{minipage}
}
\paragraph{Risposta (b)}
\begin{itemize}
\item
  Tra $V_-$ e $V_{OUT}$ i rami sono corto-circuiti, di conseguenza $V_- = V_{OUT}$.
\item
  Applico il principio di \textit{cortocircuito virtuale} e di conseguenza $V_+ = V_-$.
\end{itemize}
di conseguenza
\begin{equation*}
  V_B = V_{OUT}
\end{equation*}
le successive equazioni nodali appaiate con il risultato precedente permettono di
ottenere la funzione di trasferimento corretta
\begin{equation*}
  \begin{gathered}
    sC(V_A - V_B) = \frac{V_{OUT}}{2R} \iff V_A = \left(\frac{1 + 2sRC}{2sRC}\right)V_{OUT} \\
    sC(V_S - V_A) = sC(V_A - V_B) + \frac{V_A - V_{OUT}}{R} \\
    sRC(V_S - V_A) = sRC(V_A - V_B) + V_A - V_{OUT} \\
    sRCV_S = (1 + 2sRC)V_A - (1 + sRC)V_{OUT} \\
    sRCV_S = \left[(1 + 2sRC)\left(\frac{1 + 2sRC}{2sRC}\right) - (1 + sRC)\right]V_{OUT} \\
    2s^2R^2C^2V_S = \left[(1 + 2sRC)(1 + 2sRC) - 2sRC(1 + sRC)\right]V_{OUT} \\
  \end{gathered}
\end{equation*}
da cui
\begin{equation}
  \label{eq:hp-vcvs-butterworth-tfunc}
  \colorbox{blu}{\boxed{H(s) = \frac{2s^2R^2C^2}{1 + 2RCs + 2R^2C^2s^2}}}
\end{equation}

\noindent\colorbox{orange}{
  \begin{minipage}{1.0\linewidth}
    \paragraph{Richiesta (c)}
    Descrivere le propriet\'a analitiche di $A_v(s)$ (ovvero calcolarne zeri e poli) e da queste
    dedurre la stabilit\'a del filtro.
  \end{minipage}
}
\paragraph{Risposta (c)}
Per studiare gli zero della funzione di trasferimento complessa si studiano
le radici del numeratore della forma canonica.
\begin{equation*}
  2s^2R^2C^2 = 0
\end{equation*}
In questo caso \'e presente uno zero del secondo ordine
\begin{equation}
  \label{eq:hp-vcvs-butterworth-zeros}
  \colorbox{blu}{\boxed{z^2 = 0}}
\end{equation}
Per studiare i poli della funzione di trasferimento complessa si studiano
le radici del denominatore della forma canonica.
\begin{equation*}
  \begin{cases}
    \xi \equiv 2RCs \\
    1 + 2\xi + 2\xi^2 = 0
  \end{cases} \implies
  \xi_{1,2} = \frac{-2 \pm \sqrt{4-4\cdot 1\cdot 2}}{2\cdot 2} = -\frac12 \pm j\frac12
\end{equation*}
In questo sono presenti due poli del secondo ordine
\begin{equation}
  \label{eq:hp-vcvs-butterworth-poles}
  \colorbox{blu}{\boxed{p_{1,2} = -\frac{1}{4RC} \pm j\frac{1}{4RC}}}
\end{equation}
il filtro \'e \textbf{stabile} perch\'e entrambi i suoi poli hanno parte reale negativa. \\

\noindent\colorbox{orange}{
  \begin{minipage}{1.0\linewidth}
    \paragraph{Richiesta (d)}
    Dalla funzione di trasferimento determinare i valori attesi per le caratteristiche del filtro
    (guadagno massimo, frequenza di taglio, andamento di modulo/fase in funzione della
    frequenza) e confrontarli con quanto misurato al punto 1c.
  \end{minipage}
}
\paragraph{Risposta (d)}
Il guadagno si ottiene mediante \ref{eq:da-tfunc-a-gain}, calcolando le parti reali ed immaginarie del
numeratore e denominatore della forma canonica della funzione di trasferimento $H(s)$ si ottiene
\begin{equation*}
  \begin{gathered}
    \begin{cases}
      \mathcal{R}_{\omega}^N = -2R^2C^2\omega^2 \\
      \mathcal{I}_{\omega}^N = 0 \\
      \mathcal{R}_{\omega}^D = 1-2R^2C^2\omega^2 \\
      \mathcal{I}_{\omega}^D = 2RC\omega
    \end{cases} \implies
    G(\omega) = \sqrt{\frac{(-2R^2C^2\omega^2)^2}{(1-2R^2C^2\omega^2)^2 + (2RC\omega)^2}}
  \end{gathered}
\end{equation*}
semplificando l'espressione e sostituendo $\omega \to 2\pi f$ si ottiene il guadagno in funzione
della frequenza
\begin{equation}
  \label{eq:hp-vcvs-butterworth-gain}
  \colorbox{blu}{\boxed{G(f) = \frac{8\pi^2R^2C^2f^2}{\sqrt{1+64\pi^4R^4C^4f^4}}}}
\end{equation}
Per calcolare l'espressione della fase bisogna prima identificare il semi-piano
di appartenenza, il segno della parte reale \'e determinata da
\begin{equation*}
  (-2R^2C^2\omega^2)(1-2R^2C^2\omega^2)
\end{equation*}
il segno non \'e ben determinato per ogni $\omega \in [0, +\infty)$. Invece, il segno
della parte immaginaria \'e determinata da
\begin{equation*}
  -(-2R^2C^2\omega^2)(2RC\omega) = 4R^3C^3\omega^3 > 0
\end{equation*}
di conseguenza il semipiano \'e quello \textit{immaginario positivo} da cui
la fase $\phi$ in funzione della frequenza viene calcolata attraverso la \ref{eq:da-tfunc-a-fase-immaginario-positivo}
e la sostituzione $\omega \to 2\pi f$
\begin{equation}
  \label{eq:hp-vcvs-butterworth-fase}
  \colorbox{blu}{\boxed{\phi(f) = \frac{\pi}{2} + \arctan\left(\frac{1-8\pi^2R^2C^2f^2}{4\pi RCf}\right)}}
\end{equation}
Il numeratore e denominatore del guadagno sono di segno concorde (naturalmente), il numeratore
\'e monotono cresecente e il denominatore monotono decrescente, di conseguenza il guadagno \'e
monotono crescente di conseguenza il massimo si trova nel limite superiore del dominio
\begin{equation}
  \label{eq:hp-vcvs-butterworth-gain-max}
  \colorbox{blu}{\boxed{\max_{f \in [0, +\infty)}G(f) = \lim_{f\to +\infty}G(f) = 1}}
\end{equation}
la frequenza di taglio \'e per definizione la frequenza per cui il guadagno quadro \'e
met\'a del massimo (appena trovato)
\begin{equation*}
  \frac{8\pi^2R^2C^2f^2}{\sqrt{1+64\pi^4R^4C^4f^4}} = \frac{1}{\sqrt2} \implies
  2(8\pi^2R^2C^2f^2)^2 = 1+64\pi^4R^4C^4f^4
\end{equation*}
\begin{equation}
  \label{eq:hp-vcvs-butterworth-fcut}
  \colorbox{blu}{\boxed{f = \frac{1}{2\sqrt{2}\pi RC}}}
\end{equation}

\paragraph{Richiesta (e)}
Dare un’interpretazione qualitativa (e se possibile quantitativa) delle caratteristiche del
segnale di uscita osservato ai punti 1d, 1e.
